%!TeX root=csp-audit.tex
\section{What this document is}\label{sec:audit-intro}

The proof paper is written in the Jordan--Schönflies mould: where
\cite{ZhukSimplified} is wrong, incomplete, or ambiguous, the proof paper
simply states the corrected thing and proves it, with no marginal notes about
what the source has. That is the right shape for a document whose job is to be
checked as mathematics. It is the wrong shape for a reader who wants to know
\emph{where our rendering differs from Zhuk's paper and why}, and it would be
dishonest to make those differences invisible.

This document is the difference. For every place where the proof paper departs
from \cite{ZhukSimplified}, there is an entry here that says what the source
has, quotes it, says why it cannot be transcribed, gives the repair, and names
the statement of the proof paper that carries it.

None of this is a criticism of \cite{ZhukSimplified}, which is unusually careful
prose --- across nine thousand lines of \TeX{} it says ``obvious'' three times
and ``clearly'' never, and a mechanical sweep of \S5 turns up four hedges in
thirteen printed pages, all four of which turn out to be honest
(\S\ref{sec:hedges}). Sixteen substantive defects in fifty pages is roughly
the rate one should expect of any informal exposition of a proof this size, and
finding them is the point of reading it this way.

\subsection{The evidentiary standard}\label{sec:standard}

A claim counts as checked here only when the cited lines \emph{and the prose
around them} have been read, and the passage that does or does not discharge
the claim can be quoted. The standard is not pedantry: two of the original
claims died to it. One was taken from a reader who worked from a pseudocode
block's \texttt{Input:} line while the hypothesis it needed was in the
paragraph above (\S\ref{sec:refuted-depth}); the other asserted a citation
cycle that a parser refuted in minutes (\S\ref{sec:refuted-cycle}). Both
failure modes are the same: reading the formal-looking part of a paper and not
the sentence before it.

Where a claim is reducible to a finite check, it has had one. The searches are
listed in \S\ref{sec:machine-checks} with their scope, so that a reader can see
what they do and do not cover.

\subsection{Line numbers, and how entries are keyed}\label{sec:keying}

Line numbers are into \texttt{main.tex} and \texttt{StrongSubalgebras.tex} of
the v2 arXiv source of \cite{ZhukSimplified}, obtainable from
\texttt{arxiv.org/e-print/2404.01080v2}. Zhuk's paper is not redistributed
here, and its statement names are its own: \texttt{LEMPropagation},
\texttt{THMMainStableIntersection} and so on are the source's labels, not ours.

Each entry names the statement of the proof paper that carries the repair, by
that statement's label, set in typewriter: \pkey{lem:propagation},
\pkey{thm:stable-intersection}. The two documents are compiled separately and
share no cross-reference machinery, deliberately, so that neither can be read
as evidence for the other. The label is the join key.

The tags $\mathrm D2$, $\mathrm D4,\dots$ are historical: they were assigned in
the order the claims were raised, and some early tags were retired when the
claims they named were refuted or restated. The register of
\S\ref{sec:register} is the current and complete list; the retired claims are
in \S\ref{sec:refuted}.

\subsection{Where things stand}\label{sec:audit-tally}

\begin{center}\small
\begin{tabular}{@{}p{0.52\textwidth}r@{}}
\hline
\textbf{} & \textbf{Count} \\
\hline
Defects confirmed and repaired, needing no new mathematics & 15 \\
\quad of which: a citation the later paper dropped & 3 \\
\quad a typographical or typing slip & 2 \\
\quad a statement that is false as printed & 2 \\
\quad omitted steps and dropped hypotheses & 8 \\
Defects confirmed and \textbf{open} & 1 \\
Claims examined and refuted & 2 \\
Claims examined and declared moot & 1 \\
Readings legislated (of which two are theorems, not choices) & 10 \\
Hedges in \S5 of the source, all discharged & 4 \\
\hline
\end{tabular}
\end{center}

The single open item is $\mathrm{D}17$ (\S\ref{sec:D17}), the maximal
multi-type extension. It was declared solved in an earlier draft; an external
review found the repair invalid at its first line, and the finding is correct.
It is recorded here as open rather than quietly weakened, and the proof paper
marks the corresponding statement as open in its status table.
